\documentclass{article}
\usepackage[T1]{fontenc}
\usepackage{lmodern}
\usepackage{graphicx}
\usepackage{amsmath,amssymb}
\usepackage[a4paper,margin=2cm]{geometry}
\usepackage[absolute,overlay]{textpos}
\setlength{\TPHorizModule}{1mm}
\setlength{\TPVertModule}{1mm}
\usepackage[indonesian]{babel}
\usepackage{hyperref}
\setlength{\emergencystretch}{2em}
% unit: Halaman 82

\begin{document}

% === Halaman 82 (llm) ===

\section*{Kriptanalisis Playfair Cipher}

\begin{itemize}
    \item Karena ada 26 huruf abjad, maka terdapat $26 \times 26 = 677$ bigram, sehingga identifikasi bigram individual lebih sukar.
    \item Sayangnya ukuran poligram di dalam \textit{Playfair cipher} tidak cukup besar, hanya dua huruf sehingga \textit{Playfair cipher} tetap tidak aman.
    \item Meskipun \textit{Playfair cipher} sulit dipecahkan dengan analisis frekuensi relatif huruf-huruf, namun ia dapat dipecahkan dengan analisis frekuensi pasangan huruf.
    \item Dalam Bahasa Inggris kita bisa mempunyai frekuensi kemunculan pasangan huruf, misalnya pasangan huruf TH dan HE paling sering muncul.
\end{itemize}

\end{document}
